\section{Un texto de varios enunciados}\label{sec:texto-integrado}
% Texto original de esta revisión. Una sola leyenda para seis enunciados.
% La clave de lectura y el registro de relectura están en anexos/comprobaciones.tex.
\newcommand{\KalaContextoTexto}{%
  Documento M01: inspección ficticia de dos sensores. Los referentes se mantienen
  en todas sus páginas: 1 es el sensor A; 2 es el sensor B, distinto de A;
  3 es el registro de inspección; 4 es el evento de inspección, intervalo
  \([10,12]\); 5 es el evento de archivo, intervalo \([13,14]\).
  El tiempo se mide en segundos desde un mismo origen del experimento.
  C3 contiene exactamente esos dos sensores. P1 significa rojo según la
  clasificación declarada por el protocolo; P2 significa activo en la
  inspección. Ambas propiedades se evalúan en \(t_0=11\), bajo las mismas
  condiciones de observación. Las afirmaciones atribuidas al registro no se
  convierten por sí solas en hechos certificados.
}
\newcommand{\KalaEnunciadoUno}{\KalaFuente{\KalaReferencia{3}}{\KalaPosee{\KalaReferencia{1}}{2}}}
\newcommand{\KalaEnunciadoDos}{\KalaFuente{\KalaReferencia{3}}{\KalaNegacion{\KalaPosee{\KalaReferencia{2}}{2}}}}
\newcommand{\KalaEnunciadoTres}{\KalaFuente{\KalaReferencia{3}}{\KalaConjuncion{\KalaPosee{\KalaReferencia{1}}{1}}{\KalaNegacion{\KalaPosee{\KalaReferencia{2}}{1}}}}}
\newcommand{\KalaEnunciadoCuatro}{\KalaDefinicion{F}{\KalaSeleccion{x}{\KalaConcepto{3}}{\KalaConjuncion{\KalaPosee{\KalaVariable{x}}{1}}{\KalaPosee{\KalaVariable{x}}{2}}}}}
\newcommand{\KalaEnunciadoCinco}{\KalaPregunta{\KalaExiste{x}{\KalaMat{F}}{\KalaMiembro{\KalaVariable{x}}{\KalaConcepto{3}}}}}
\newcommand{\KalaEnunciadoSeis}{\KalaRelacion{1}{\KalaReferencia{4}\KalaSeparador\KalaReferencia{5}}}
\newcommand{\KalaMostrarEnunciado}[2]{%
  \par\noindent{\small\sffamily\color{KalaSuave}#1}\par
  \KalaMostrar{#2}\par\vspace{2mm}%
}

Con este ejemplo exploro cómo mantener referencias y contexto en varios enunciados.
\par\KalaContextoTexto
\par\textbf{Convención de página:} M01 identifica todo el documento; U1--U6
son números editoriales de enunciado. No son nuevos signos de KalaScript.
El nombre local \(F\) se introduce en U4 y conserva esa definición en U5.

\KalaMostrarEnunciado{M01 / U1}{\KalaEnunciadoUno}
\KalaMostrarEnunciado{M01 / U2}{\KalaEnunciadoDos}
\KalaMostrarEnunciado{M01 / U3}{\KalaEnunciadoTres}
\clearpage
\noindent{\small\sffamily M01 · continuación; se conserva la leyenda anterior.}
\KalaMostrarEnunciado{M01 / U4}{\KalaEnunciadoCuatro}
\KalaMostrarEnunciado{M01 / U5}{\KalaEnunciadoCinco}
\KalaMostrarEnunciado{M01 / U6}{\KalaEnunciadoSeis}

\subsection{Qué permite comprobar este texto}
El mismo referente 3 se reutiliza en tres atribuciones; los referentes no se
reinician al pasar de página. La negación de U2 afecta al estado activo de B,
mientras que en U3 afecta solo a su propiedad de color. U4 define un conjunto
y U5 formula una pregunta sobre él. U6 establece un orden temporal con eventos
diferentes de los sensores. Conservo la clave detallada y mis notas de
relectura en el anexo de comprobaciones del manual.
